% RESEARCH_DRAFT_SHA256: 76475a46f16e42f7d0910bec1d911392f859a743c8a7729e6a5343a3b1d627c7
% Layout and source-reference wrappers expanded for the web renderer.
\documentclass[11pt]{article}
\usepackage[margin=25mm]{geometry}
\usepackage{amsthm}
\usepackage{fontspec,unicode-math}
\setmainfont{Latin Modern Roman}
\setmathfont{Latin Modern Math}
\usepackage{amsmath,amssymb,mathtools}
\usepackage{xurl,hyperref}
\hypersetup{colorlinks=true,urlcolor=blue,linkcolor=blue}
\setcounter{secnumdepth}{0}
\setlength{\parindent}{0pt}
\setlength{\parskip}{5pt}
\setlength{\emergencystretch}{3em}
\newcommand{\sref}[2]{\hyperlink{source:#1:#2}{[S#2]}}
\newcommand{\cataloguescope}[1]{\paragraph{Source question.}#1}
\newcommand{\floor}[1]{\left\lfloor#1\right\rfloor}
\newtheorem{proposition}{Proposition}
\newtheorem{lemma}[proposition]{Lemma}
\begin{document}
\paragraph{Research manuscript.} OpenAI \texttt{gpt-6-astra}, reasoning \texttt{ultra}; 11 October 2026.
\paragraph{Preserved source body.} The complete catalogue mathematical text and sources follow with layout and source links expanded, including source qualifications and its historical preparation notes. New research begins at the section headed Research attempt.

\section{Covariance-whitened polynomial invariants for model goodness of fit}\label{statement}

\subsection{Definitions and mathematical statement}

\paragraph{Shared notation.}Standard number systems are denoted by $\mathbb N,\mathbb Z,\mathbb Q,\mathbb R,\mathbb C$; index ranges, computational models and specialized notation follow the cited problem source and the definitions below.


A model invariant is a polynomial $f$ vanishing on every probability vector in the model. For observed counts with total $n$, the empirical probability vector is $\widehat p_i=x_i/n$. Evaluating invariants at $\widehat p$ gives residuals; aggregation such as $\sum_f|f(\widehat p)|$ depends on the chosen invariant set and its scaling.
\paragraph{Statement recorded in the supplied source.}
Question 61 (From graphical models notes). Is there any useful statistical measure that can be obtained from evaluating the polynomials that vanish on probabilities coming from a given model at an empirical distribution? What about

∑ |f (̂p)|? \href{https://aimath.org/WWN/compalgstat/compalgstat.pdf}{[S2]}
\paragraph{Scope qualification.}
The supplied extraction includes adjacent workshop exposition; the numbered target passage is retained verbatim and that surrounding exposition is omitted. No normalization of invariants, null distribution, or loss criterion is supplied; an unnormalized residual sum is not claimed to be a calibrated test.
\paragraph{Source status.}
The supplied catalogue labels this record partially solved. Its source-status metadata was checked against the supplied snapshot on 2026-10-11. This import does not independently establish current open status or mathematical correctness.
\subsection{Short English statement}
Can polynomial residuals evaluated on observed frequencies provide a useful statistical measure of model fit?
\subsection{Sources}
\begin{itemize}
\item[S1] Ulam AI and the UnsolvedMath Contributors, UnsolvedMath: A Curated Collection of Open Mathematics Problems, catalogue entry. Catalogue curation and metadata: CC BY 4.0. The supplied source question is retained; definitions and the English summary are editorial additions.
\url{https://huggingface.co/datasets/ulamai/UnsolvedMath}.
\item[S2] AIM Workshop Problem Lists, original problem source linked by the catalogue
\url{https://aimath.org/WWN/compalgstat/compalgstat.pdf}.
\end{itemize}
\paragraph{Preparation and review.}
Original question attributed to the cited source and UnsolvedMath contributors. Codex (GPT-6; exact serving revision not exposed) assisted with definitions, English summaries, source inspection and formatting. No new human mathematical authorship or independent proof verification is asserted. Catalogue identity checked against the complete supplied registry and existing repository statements; known flagged repetitions and obvious identical conjectures excluded. Full mathematical-equivalence/literature review has not been performed. Statement preparation and final review followed contribution prompts C and D; automated checks verify structure and identity, not open status.
\label{end:statement}


\section{Research attempt}
\subsection{Literature and scope}
Question 61 suggests scoring an empirical distribution by values of vanishing polynomials, such as a sum of absolute residuals. I tested invariance under replacing a polynomial by another with the same zero set. Powers expose an obstruction even after one recognizes ordinary scalar-rescaling dependence. This elementary example is not a new statistical test.
\begin{proposition}
The absolute value of an arbitrary defining polynomial, even one with the same zero set as a regular independence equation, cannot by itself give a representation-independent measure of model discrepancy. First-order covariance normalization also fails when that defining polynomial is a nontrivial power.
\end{proposition}
\begin{proof}
For a binary joint probability table set $g(p)=p_{00}p_{11}-p_{01}p_{10}$. For every integer $r\ge1$, the polynomial $f_r=g^r$ has the same real and complex zero set as $g$, and vanishes on the independence model. At the fixed positive empirical table
\[
\widehat p=\begin{pmatrix}7/20&3/20\\3/20&7/20\end{pmatrix}
\]
one has $g(\widehat p)=1/10$, so $|f_r(\widehat p)|=10^{-r}\longrightarrow0$. The empirical distribution has not moved and remains outside the model. Thus the proposed residual can be made arbitrarily small solely by changing its defining polynomial.

For $r\ge2$, $\nabla f_r=r g^{r-1}\nabla g$ vanishes at every model point. In contrast, at the uniform table the gradient of $g$, in order $00,01,10,11$, is $(1,-1,-1,1)/4$, which is nonzero and tangent to the probability-sum hyperplane. A first-order variance expression of the form $\nabla f_r^{\mathsf T}\Sigma\nabla f_r$ is therefore zero there for any covariance matrix $\Sigma$, so dividing by its square root is undefined. This degeneracy arises from the chosen polynomial, even though the original equation is regular on that hyperplane at this point.
\end{proof}
\subsection{Verification and remaining gaps}
Rational checks gave residuals $1/10,1/100,1/100000,1/10000000000$ for powers $1,2,5,10$. Powers generate a nonradical ideal when used alone; they are not asserted to generate the same ideal, only the same zero set and valid vanishing polynomials. A method using a prescribed regular generating system and suitable normalization can avoid this particular obstruction. Constructing useful, calibrated tests, including singular cases and generator selection, remains open within this attempt.

\section{References and source attribution}
The original question is credited to the AIM workshop list prepared by Seth Sullivant (2004), \url{https://aimath.org/WWN/compalgstat/compalgstat.pdf}. The source question and catalogue attribution are reproduced above. This list was accessed on 11 October 2026; it is historical evidence of the question, not a certification of present open status.


\paragraph{Integration review.} Codex (GPT-6; exact serving revision not exposed) reviewed the written arguments and expanded layout and source-reference wrappers. Repository metadata and the production web renderer were checked using Python 3.13 and JavaScript with MathJax 3.2.2. This model-assisted review is not independent human proof verification; no PDF compilation or formal proof checking is asserted.

\section{Completion Estimate}
\noindent\textbf{COMPLETION: 0\%}
No substantiated new progress toward the entire original target is claimed. The worked results and failed route are delimited above; a known result, illustrative model, or auxiliary obstruction is not counted as original-target completion.

\section{Final status and attribution}
\textbf{Attempt status: unresolved.} This describes this research attempt, not a new ruling on the source catalogue's status or on external literature claims.
The mathematics and this self-audit were generated by OpenAI \texttt{gpt-6-astra}, reasoning \texttt{ultra}, on 11 October 2026. Codex (GPT-6; exact serving revision not exposed) orchestrated the batch and contributed formatting/repository review. No human mathematical authorship, independent human proof checking, or formal verification is claimed. The historical preparation paragraphs in the reproduced catalogue body concern that body only.

Contribution prompt A was applied to the complete definition and supplied source record, followed by the adversarial checks of prompt B. Local reproduction files are \texttt{batch005/checks.py}, \texttt{computations.json}, and the per-ID audit JSON. They were not public links at preparation time. Numerical checks are finite and do not prove unrestricted results. No LaTeX engine was found locally; compilation was not performed by this research agent.

\end{document}
